\scanpage{30}
\mspage{29}

l'identité). Il transforme $0,1,\infty$ en
\[
7)\quad\left\{\begin{array}{l}\sigma0=2\\\sigma1=-1\\\sigma\infty=\tfrac12\end{array}\right.
\]
\margin{Prouver : le normalisateur de $\mathfrak S_3$ dans $\operatorname{Aut}P^1_S$ est $\mathfrak S_3\times(\unsure{\mathbb Z}/2)$}

La différente de $\mathfrak S_3$ -- ou ens pts de ramification -- est \unsure{formé} des ens. des pts \unsure{de} un des cinq opérateurs $\lambda\mu\nu\varphi\psi$ \unsure{distincts} de l'identité. Pour $\lambda,\mu,\nu$ on trouve les sections $(\unsure{1},-1)(\infty,\tfrac12)(\unsure{\infty|0},2)$ ; [\unsure{st.ble} par $\sigma$]. Pour $\varphi$ (et $\psi$, c'est pareil), on trouve $1-\frac1t=t$ ou
\[
8)\qquad t^2-t+1=0
\]
i.e. l'équation cyclotomique
\[
\boxed{\Phi_6(t)=0}.
\]
\deleted{\unsure{Le}}

\deleted{\unsure{distinct de $\sigma$}}

C'est aussi l'\unsure{ens.} $\overline S$ des pts fixes de $\sigma$.

[Il est non ramifié \inserted{sur $S$} aux pts $s\in S$ tels que $k(s)$ soit de car. $\ne2,3$.].
\[
\mathbb Z_2=\mathfrak S_3/(\mathbb Z/3)
\]
y opère, de façon \uline{non triviale} aux pts où il est non ramifié sur $S$.

D'autre part, le schéma \inserted{$G$} des automorphismes de $P^1$ qui laissent $\overline S$ \unsure{globalement} invariant est une extension de $\mathbb Z_2$ [automorphismes de $\overline S$] par $(G_m)_S$ \uline{tordu} ---- [par le rev. cyclotomique]
\margin{\unsure{les except\textsuperscript{ns} qu'il y a pour car. $2$ ou $3$ de $S$}\quad\figure{\begin{tikzpicture}\draw (0,0)--(0,1.6);\end{tikzpicture}}}
